
\chapter{Long, wide, and ``tidy'' data}

\section{Wide form data}

Consider the following \texttt{DataFrame} (called ``\texttt{g}''), which
contains a professor's grade book for his students:

\begin{Verbatim}[fontsize=\small,samepage=true,frame=single,framesep=3mm]
print(g)
\end{Verbatim}
\vspace{-.2in}

\begin{Verbatim}[fontsize=\small,samepage=true,frame=leftline,framesep=5mm,framerule=1mm]
       name year    Q1    Q2    Q3    Q4    Q5    Q6    Q7    Q8
0   Beyonce   So   9.0  10.0   7.0  10.0   8.0  10.0   9.0  10.0
1     Jay-Z   Jr   7.0  10.0   4.0  10.0  10.0   8.0   3.0  10.0
2    Bortis   Sr   8.0  10.0   7.0   3.0   7.0   7.0   7.0   9.0
3       Han   So   NaN   1.0   NaN  10.0   7.0   8.0   7.0   2.0
4       Mal   Jr   5.0   4.0   9.0   NaN   7.0   2.0   5.0  10.0
5     Inara   Jr   3.0   3.0   4.0  10.0   6.0   4.0   7.0   5.0
6   Obi-wan   So   5.0  10.0   9.0   4.0   1.0   4.0   NaN   2.0
7      Finn   Jr   6.0   3.0   5.0   6.0  10.0   NaN   NaN   3.0
\end{Verbatim}

It's not hard to tell what's going on: each student's name and year in school
is recorded, as is their grade on the eight class quizzes. Some students took
zeroes on one of more quizzes (like \texttt{Han} did for quizzes 1 and 3), and
the quiz grades range from 0 to 10.

\index{wide form}

The structural style for this kind of \texttt{DataFrame} is called \textbf{wide
form}. Each piece of a student's data is in its own column, making the
\texttt{DataFrame} quite ``wide,'' and comparatively not very long (tall).

Now there's nothing wrong with wide form. All the data is there, and it allows
us to ask many questions easily, like:

\begin{compactitem}
\item[\leftthumbsup] What was Inara's score on Quiz \#5? \ (\texttt{g[g.name=="Inara"].Q5})
\item[\leftthumbsup] What was the high score on Quiz \#3? \ (\texttt{g.Q3.max()})
\item[\leftthumbsup] What was Mal's average quiz score? \ 
(\texttt{g[g.name=="Mal"].mean(axis=1)})
\end{compactitem}

\smallskip
Other kinds of questions, though, defy easy answers, like these:

\begin{compactitem}
\item[\leftthumbsdown] What was the overall average score (for all quizzes)? \ (?)
\item[\leftthumbsdown] How did class performance compare between the first half
of the semester (quizzes 1-4) and the second half (quizzes 5-8)? \ (?)
\item[\leftthumbsdown] How did freshmen, sophomores, juniors, and seniors
compare? \ (?)
\end{compactitem}

The reason these questions are difficult are that we have quiz scores spread
across different columns. As long as we're happy treating them in groups like
that, it's easy, but as soon as we want to carve up the quiz scores in some
other way, it becomes hard.

\section{Long form data}

\index{melt@\texttt{melt()} (Pandas)}

The solution is to convert the \texttt{DataFrame} to a structure which
collapses all the quiz data into a single column. For this, we use the
\texttt{melt()} function from the Pandas package:

\begin{Verbatim}[fontsize=\small,samepage=true,frame=single,framesep=3mm]
gl = pd.melt(g, ["name","year"], var_name="quiz", value_name="score")
\end{Verbatim}

Now this is kind of complicated so pay attention. 

\index{list}

\begin{compactitem}

\item The first argument is the name of the wide-form \texttt{DataFrame} you
want to ``melt.''

\item The second argument is a list of the columns \textit{that you want to
leave alone.} In other words, a list of ``off-limits'' columns that you want
Pandas \textit{not} to melt.

\item The third argument, named \texttt{var\_name}, is \textit{the name you
want for your new column whose values will be the old column
\underline{names}.} It can be hard to get your head around this one. You have
to imagine all the current column names that you're melting being put in a
column, and ask yourself: ``what is the right name for that new column?'' In
this case, since we're melting \texttt{Q1}, \texttt{Q2}, \textit{etc.}, the
right name for the column is something like ``\texttt{quiz},'' since those are
the quiz numbers.

\item Finally, the fourth argument, named \texttt{value\_name}, is \textit{the
name you want for your new column whose values will be the old column's
\underline{values}.} You have to ask yourself: ``what are those 'things' that
are currently values in the \texttt{Q1}, \texttt{Q2}, \textit{etc.} columns?''
A reasonable answer would be ``\texttt{score},'' since each one is a quiz
score.

\end{compactitem}

We set the result of the \texttt{melt()} call to a variable called \texttt{gl}
(where ``\texttt{l}'' stands for ``long.'') Here's what it looks like:

\begin{Verbatim}[fontsize=\small,samepage=true,frame=single,framesep=3mm]
print(gl)
\end{Verbatim}
\vspace{-.2in}

\begin{Verbatim}[fontsize=\small,samepage=true,frame=leftline,framesep=5mm,framerule=1mm,xleftmargin=2cm]
       name year quiz  score
0   Beyonce   So   Q1    9.0
1     Jay-Z   Jr   Q1    7.0
2    Bortis   Sr   Q1    8.0
3       Han   So   Q1    NaN
4       Mal   Jr   Q1    5.0
5     Inara   Jr   Q1    3.0
6   Obi-wan   So   Q1    5.0
7      Finn   Jr   Q1    6.0
8   Beyonce   So   Q2   10.0
9     Jay-Z   Jr   Q2   10.0
10   Bortis   Sr   Q2   10.0
11      Han   So   Q2    1.0
12      Mal   Jr   Q2    4.0
13    Inara   Jr   Q2    3.0
14  Obi-wan   So   Q2   10.0
15     Finn   Jr   Q2    3.0
16  Beyonce   So   Q3    7.0
17    Jay-Z   Jr   Q3    4.0
18   Bortis   Sr   Q3    7.0
19      Han   So   Q3    NaN
20      Mal   Jr   Q3    9.0
21    Inara   Jr   Q3    4.0
22  Obi-wan   So   Q3    9.0
23     Finn   Jr   Q3    5.0
            . . .
60      Mal   Jr   Q8   10.0
61    Inara   Jr   Q8    5.0
62  Obi-wan   So   Q8    2.0
63     Finn   Jr   Q8    3.0
\end{Verbatim}

\index{long form}

This is called \textbf{long form}, for obvious reasons. We now have one row for
every single quiz score. It might seem wordy, and I'll admit it's not as
easy to locate and eyeball a single result. However, all our ``hard'' questions
from above can be easily answered:

\begin{compactitem}
\item[\leftthumbsup] What was the overall average score (for all quizzes)? \
(\texttt{gl.score.mean()})
\item[\leftthumbsup] How did class performance compare between the first half
of the semester (quizzes 1-4) and the second half (quizzes 5-8)? \\
(\texttt{gl[gl.quiz.isin(["Q1","Q2","Q3","Q4"])].mean()})
\item[\leftthumbsup] How did freshmen, sophomores, juniors, and seniors
compare? \\ (\texttt{gl.groupby("year").score.mean()})
\end{compactitem}

Turning what were column headings into \textit{values} gave us all the
flexibility we needed. We can still treat all the scores individually by quiz
if we want to, but we can also treat them as a whole, or aggregate them in
other ways (like by school year).

As an aside, one thing I'd probably do to this \texttt{DataFrame} is get rid of
the leading ``\texttt{Q}'' before each \texttt{quiz} value, and convert that
column to integers to boot:

\begin{Verbatim}[fontsize=\footnotesize,samepage=true,frame=single,framesep=3mm,xleftmargin=1cm]
gl.quiz = gl.quiz.str[1]
gl.quiz = gl.quiz.astype(int)
print(gl)
\end{Verbatim}

\vspace{-.2in}

\begin{Verbatim}[fontsize=\scriptsize,samepage=true,frame=leftline,framesep=5mm,framerule=1mm,xleftmargin=1cm]
       name year  quiz  score
0   Beyonce   So     1    9.0
1     Jay-Z   Jr     1    7.0
2    Bortis   Sr     1    8.0
3       Han   So     1    NaN
4       Mal   Jr     1    5.0
5     Inara   Jr     1    3.0
6   Obi-wan   So     1    5.0
7      Finn   Jr     1    6.0
8   Beyonce   So     2   10.0
9     Jay-Z   Jr     2   10.0
10   Bortis   Sr     2   10.0
11      Han   So     2    1.0
12      Mal   Jr     2    4.0
13    Inara   Jr     2    3.0
14  Obi-wan   So     2   10.0
15     Finn   Jr     2    3.0
16  Beyonce   So     3    7.0
17    Jay-Z   Jr     3    4.0
18   Bortis   Sr     3    7.0
\end{Verbatim}

\index{str suffix@\texttt{.str} suffix (Pandas)}

The ``\texttt{.str}'' suffix, you'll recall, lets us perform standard string
operations (like \texttt{len()}, \texttt{contains()}, or indexing) on entire
\texttt{DataFrame} columns. Here, by saying ``\texttt{gl.quiz.str[1]}'', we're
requesting only the \textit{second} character (character \#1) of each string,
which extracts the ``\texttt{1}'' from ``\texttt{Q1}'' and so forth.

Treating quiz numbers like the integers they are is conceptually correct, and
also allows us to do things like query for things like ``quizzes after quiz
\#3'' instead of having to list a bunch of categorical values.

\section{``Tidy'' data}

\index{Wickham, Hadley}

The term ``tidy'' to describe a data set has sprung into prominence recently
as a result of the excellent paper by Hadley Wickham (of R fame) who neatly
captured the perils of various data formatting choices. Hadley first defines
these terms:

\begin{compactitem}
\item \textbf{variable} -- a category of something, for which individual
objects have a value. Examples: name, height, gender, center field distance.
\item \textbf{value} -- an actual individual piece of data corresponding to a
variable for one object. Examples: Kylo Ren, 5'6", female, 390 feet.
\item \textbf{observation} -- a particular measurement that comprises all the
information necessary to interpret that measurement. Examples: Jeff's
weight was 165 lbs., Jenny's golf score at Field \& Cross Links in the rain
was 88. (In the first example, ``all the information necessary'' includes
``Jeff'' and ``165''; in the second, it includes ``Jenny'', ``Field \& Cross
Links'', ``rainy'', and ``88''.)
\item \textbf{observational unit} -- a collection of all the observations of a
particular type. (``Golf performances'' would be an observational unit,
separate from a ``golfers'' unit which contains information about golfers
themselves, not their particular performances.)
\end{compactitem}

\index{tidy@``tidy'' data}

\vspace{.2in}

Then, he outlined one simple rule. A data set is considered \textbf{tidy} if
and only if:

\definecolor{shadecolor}{rgb}{.9,.9,.9}
\begin{shaded}
\begin{compactenum}
\item \label{ownTable} Each observational unit is in its own table.
\item \label{obsRow} In each table, each observation forms a row.
\item \label{colHeader} In each table, each column header is a variable, and
the column contains values.
\end{compactenum}
\end{shaded}

\index{normalized (data set)}

Point \#\ref{ownTable} is essentially the same as ``normalized,'' from
section~\ref{normalized} (p.\pageref{normalized}). Points \#\ref{obsRow} and
\#\ref{colHeader} are essentially ``long form,'' as described in the previous
section. Hence, I usually think of ``tidy'' as meaning ``normalized and long.''

Our \texttt{gl} grade book was long, but not tidy, because it had repeated
information in it: namely, every student's \texttt{year} was repeated for every
one of their quiz scores. To normalize it, we'd want to divide it into two
tables; say, \texttt{student\_years} and \texttt{grades}:

\index{double boxies@``double boxies''}
\index{drop\_duplicates@\texttt{.drop\_duplicates()} (Pandas)}
\index{drop@\texttt{.drop()} (Pandas)}
\index{axis argument@\texttt{axis} argument}

\begin{Verbatim}[fontsize=\small,samepage=true,frame=single,framesep=3mm]
student_years = gl[['name','year']].drop_duplicates()
grades = gl.drop('year',axis=1)
\end{Verbatim}

The \texttt{.drop\_duplicates()} method gets rid of identical rows, retaining
only one. The \texttt{.drop()} method keeps all columns except for one, which
is convenient if you don't want to type all the names of the ones you want to
keep. (``\texttt{axis=1}'' means ``I'm talking about columns here, not rows.'')

Our two remaining (tidy!) \texttt{DataFrames} are now:

\begin{Verbatim}[fontsize=\small,samepage=true,frame=single,framesep=3mm,xleftmargin=1cm]
print(student_years)               print(grades)
\end{Verbatim}

\begin{Verbatim}[fontsize=\small,samepage=true,frame=leftline,framesep=5mm,framerule=1mm,xleftmargin=1cm]
      name year                           name  quiz  score
0  Beyonce   So                    0   Beyonce     1    9.0
1    Jay-Z   Jr                    1     Jay-Z     1    7.0
2   Bortis   Sr                    2    Bortis     1    8.0
3      Han   So                    3       Han     1    NaN
4      Mal   Jr                    4       Mal     1    5.0
5    Inara   Jr                    5     Inara     1    3.0
6  Obi-wan   So                    6   Obi-wan     1    5.0
7     Finn   Jr                    7      Finn     1    6.0
                                   8   Beyonce     2   10.0
                                   9     Jay-Z     2   10.0
                                   10   Bortis     2   10.0
                                   11      Han     2    1.0
                                   12      Mal     2    4.0
                                   13    Inara     2    3.0
                                   14  Obi-wan     2   10.0
                                   15     Finn     2    3.0
                                   16  Beyonce     3    7.0
                                   17    Jay-Z     3    4.0
                                   18   Bortis     3    7.0
\end{Verbatim}

Every piece of information in one and only one place. Very tidy indeed.

\subsection{Messy data sets}

Some data sets don't make it easy on us. Check out this decidedly
\textit{un}tidy one, which contains survey results about income and religious
affiliation:

\small
\begin{Verbatim}[fontsize=\small,samepage=true,frame=none]
            religion   <$10k   $10-20k  $20-30k  $30-40k   $40-50k  $50-75k
0           Agnostic      27        34       60       81        76      137
1            Atheist      12        27       37       52        35       70
2           Buddhist      27        21       30       34        33       58
3           Catholic     418       617      732      670       638     1116
4  Dont know/refused      15        14       15       11        10       35
...
\end{Verbatim}
\normalsize

As before, the problem here is that \textit{the column headings are actually
values, not variables}. Think about it: ``\texttt{\$10-20k}'' isn't the name of
a variable which might have actual values, but the value (range) itself of some
variable.

Again the two questions: (1) What would the variable be for these column
values? Answer: \textit{Income.} (2) What are the cells themselves
representing? Answer: \textit{Frequency.} (\textit{i.e.}, the number of people
polled in the survey who had that religion and that income level.)

Now we can tidy this with \texttt{melt()}:

\begin{Verbatim}[fontsize=\footnotesize,samepage=true,frame=single,framesep=3mm]
tdf = pd.melt(df, ['religion'], var_name='income', value_name='frequency')
\end{Verbatim}
\vspace{-.2in}

\begin{Verbatim}[fontsize=\small,samepage=true,frame=leftline,framesep=5mm,framerule=1mm]
       religion    income  frequency
  0    Agnostic     <$10k         27
  1     Atheist     <$10k         12
  2    Buddhist     <$10k         27
  3    Catholic     <$10k        418
 ...
\end{Verbatim}

\index{pivot}

By the way, the inverse operation of ``melt'' is often called
``\textbf{pivot}''. To reverse the above transformation in Pandas, you can use
the \texttt{.pivot()} method:

\begin{Verbatim}[fontsize=\footnotesize,samepage=true,frame=single,framesep=3mm]
wide = long.pivot(index='religion', columns='income', values='frequency')
\end{Verbatim}

and get back your original wide form \texttt{DataFrame}.


% Exercise:
% load http://cs.umw.edu/~stephen/data219/students.csv and tidy it to:
%
%                      students
%            name  year major    school  creds
%   Mal A. Kenobi  Soph  ECON       UMW     32
%  Jay O. Aniston   Sen  ECON  Richmond     92
%
%               gpas
%             name year   GPA
%   Mal A. Kenobi   So  2.47
%  Jay O. Aniston   Sr  2.30
%      (no NaNs)
%
%           schools
%    school     type   mascot
%       UMW   public   Eagles
%  Richmond  private  Spiders
%
% DELETE the original table.
% now add a column for "urban/rural" to schools table. (We'll call Richmond,
% Fredericksburg, and Charlottesville urban, and the others rural.)
%
% how many students are there?
% how many total year-GPA measurements are there?
% what's the overall GPA?
% how many students in rural vs. urban schools?
% average freshman GPA in rural vs. urban schools?
% visual: do freshmen vs. non-freshmen GPAs differ? (will require another col)
% visual: does avg GPA vary by major? by school? (will require another col)
% scatterplot of creds vs. average-yearly-GPA (ditto)
%
% helpful:
% 1) after splitting out school into separate df:
%  df.drop_duplicates()
% 2) to remove NaN's from GPAs once melted:
%  np.isnan(...)   <-- a ufunc
%  ~   <-- boolean "not" operator
%  (to set up: "gpas = gpas[~np.isnan(gpas.GPA)]")
%  (could also do: "gpas.drop(gpas.index[np.isnan(gpas.GPA)])")
% 3) df.col.str   to perform string operations as ufunc
%  (to set up: "gpa.year = gpa.year.str[3:]")
%
% 
